\chapter{Worked Examples}
\label{chap:worked_examples}

This chapter contains a series of end-to-end worked examples designed to
illustrate typical \actsim workflows. Each example states its purpose,
describes the input data, presents the Python code, and discusses the
output and its actuarial interpretation.

% =====================================================================
\section{Example 1: Aggregate Loss Simulation}
\label{ex:agg_loss}

\paragraph{Purpose.} Simulate aggregate losses for a single line of
business using a Poisson frequency distribution and a lognormal severity
distribution.

\paragraph{Input data.} No external data; the example uses synthetic
distribution parameters.

\paragraph{Code.}
\begin{lstlisting}[style=pythonstyle]
from actsim import StochasticSimulator

simulator = StochasticSimulator(
    freq_dist="poisson",
    freq_params=(50,),
    sev_dist="lognormal",
    sev_params=(10, 0.5),
    num_sim=50000,
    keep_all=True,
    seed=2024,
)

simulator.gen_agg_simulations()

print("Mean aggregate loss:",  simulator.results.mean())
print("99% VaR:",              simulator.calc_agg_percentile(99))
print("99.5% VaR:",            simulator.calc_agg_percentile(99.5))
\end{lstlisting}

\paragraph{Output interpretation.} The mean aggregate loss should be
close to $E[N] \cdot E[X]$. Because the severity is lognormal with
parameters $\mu = 10$ and $\sigma = 0.5$, $E[X] \approx
\exp(\mu + \sigma^{2}/2)$. The simulated mean should converge to this
theoretical mean as the number of simulations increases.

\paragraph{Actuarial notes.} This example serves as a sanity check
that the simulator correctly implements the compound Poisson--lognormal
model. The 99\% VaR and 99.5\% VaR are typical risk indicators reported
in solvency frameworks.

% =====================================================================
\section{Example 2: Fitting Severity Distributions to Historical Losses}

\paragraph{Purpose.} Fit a set of candidate severity distributions to a
synthetic dataset and select the best-fitting distribution by AIC.

\paragraph{Input data.} A simulated dataset of 10{,}000 losses drawn
from a known lognormal distribution. In a real engagement, this dataset
would be replaced with empirical losses from claim files.

\paragraph{Code.}
\begin{lstlisting}[style=pythonstyle]
from actsim import DistributionFitter, load_config
from actstats import actuarial as act

# Generate synthetic loss data
loss_data = act.lognormal(0.5, 0.2).rvs(size=10000)

# Load the default configuration
config = load_config()
distributions = config.distributions["severity"]
metrics = config.metrics

# Fit candidate distributions
fitter = DistributionFitter(loss_data,
                            distributions=distributions,
                            metrics=metrics)
fitter.fit()

# Inspect best fits and selected distribution
print("Best fit by AIC:", fitter.best_fits["aic"]["name"])
print("Selected fit:",    fitter.selected_fit["name"])

# Diagnostic plot
fitter.plot_predictions()

# Summary table
print(fitter.summary())
\end{lstlisting}

\paragraph{Output interpretation.} The selected distribution should
typically be the lognormal, since the data were drawn from that family.
The summary table reports AIC, BIC, log-likelihood, and chi-square
values for each fitted distribution.

\paragraph{Actuarial notes.} In real-world fits, the differences between
candidate distributions in AIC and BIC are often small. Diagnostic plots
and tail-fit checks should be consulted in addition to the
information criteria.

% =====================================================================
\section{Example 3: Synthetic Policy and Claim Generation}

\paragraph{Purpose.} Generate a synthetic dataset of 100 policies and
their associated claims.

\paragraph{Input data.} A randomly generated policy table.

\paragraph{Code.}
\begin{lstlisting}[style=pythonstyle]
import pandas as pd
import numpy as np
from actsim import ClaimSimulator

policies = pd.DataFrame({
    "policy_id":   range(1, 101),
    "freq_dist":   "poisson",
    "freq_params": list(zip(np.random.uniform(0.6, 0.8, 100).round(2),)),
    "sev_dist":    "lognormal",
    "sev_params":  list(zip(
        np.random.uniform(8, 12, 100).round(2),
        np.random.uniform(0.3, 0.7, 100).round(2),
    )),
    "start_date":  pd.Timestamp("2023-01-01"),
    "end_date":    pd.Timestamp("2023-12-31"),
})

claim_sim = ClaimSimulator(policies, random_seed=42)
claim_sim.simulate_claims()
print(claim_sim.claim_data.head())
\end{lstlisting}

\paragraph{Output interpretation.} The \texttt{claim\_data} DataFrame
contains one row per simulated claim, with a \texttt{policy\_id} column
linking each claim back to its parent policy.

\paragraph{Actuarial notes.} Some policies may produce zero claims; this
is expected behaviour for low-frequency policies and is consistent with
the underlying Poisson assumption.

% =====================================================================
\section{Example 4: Developing Claims Using LDFs}

\paragraph{Purpose.} Develop the simulated claims from
Example~3 using a user-supplied LDF schedule.

\paragraph{Code.}
\begin{lstlisting}[style=pythonstyle]
claim_sim.simulate_dates_nhpp(lambda0=10, alpha=0.5, phase=0, T=1)

base_LDFs = {0: 2.0, 3: 1.5, 6: 1.2, 9: 1.1,
             12: 1.05, 15: 1.02, 18: 1.0}

claim_sim.simulate_claim_development(
    base_LDFs=base_LDFs,
    volatility=0.1,
    cumulative_factor=1.0,
)

print(claim_sim.claim_development.head())
\end{lstlisting}

\paragraph{Output interpretation.} The \texttt{claim\_development}
DataFrame contains one row per (claim, development month) pair, with
incurred losses that grow towards the ultimate value as the development
age increases.

% =====================================================================
\section{Example 5: Converting Simulated Claims to a Triangle}

\paragraph{Purpose.} Aggregate the simulated claim development data
into a long-format loss triangle.

\paragraph{Code.}
\begin{lstlisting}[style=pythonstyle]
triangle_df = (
    claim_sim.claim_development
    .groupby(["accident_year", "development_month"])["incurred_loss"]
    .sum()
    .reset_index()
)
print(triangle_df.head())
\end{lstlisting}

\paragraph{Output interpretation.} \texttt{triangle\_df} is the
long-format triangle ready for use with \texttt{chainladder}.

% =====================================================================
\section{Example 6: Analysing the Triangle with chainladder}

\paragraph{Purpose.} Fit a chain-ladder model to the simulated
triangle and obtain projected ultimate losses.

\paragraph{Code.}
\begin{lstlisting}[style=pythonstyle]
import chainladder as cl

triangle = cl.Triangle(
    triangle_df,
    origin="accident_year",
    development="development_month",
    columns="incurred_loss",
    cumulative=True,
)

model = cl.Chainladder().fit(triangle)
print(model.ultimate_)
\end{lstlisting}

\paragraph{Output interpretation.} The projected ultimates can be
compared against the sum of simulated ultimate losses by accident year.
For a well-specified simulation and a sufficiently large sample, the
two should agree within sampling error.

% =====================================================================
\section{Example 7: Calculating VaR and TVaR}

\paragraph{Purpose.} Compute VaR and TVaR at multiple confidence levels
from a simulation run.

\paragraph{Code.}
\begin{lstlisting}[style=pythonstyle]
from actsim import StochasticSimulator

simulator = StochasticSimulator(
    "poisson", (100,),
    "gamma", (2, 5000),
    num_sim=50000,
    keep_all=True,
    seed=99,
)

simulator.gen_agg_simulations()
report = simulator.analyze_results(quantiles=[0.5, 0.9, 0.95, 0.99, 0.995])
print(report.round(2))
\end{lstlisting}

\paragraph{Output interpretation.} The output table contains VaR, TVaR,
OEP, and AEP at each requested quantile. TVaR values exceed the
corresponding VaR at the same level, as expected.

% =====================================================================
\section{Example 8: Validating Simulated Frequency}

\paragraph{Purpose.} Compare the simulated mean claim count against the
theoretical expected value.

\paragraph{Code.}
\begin{lstlisting}[style=pythonstyle]
import numpy as np
from actsim import StochasticSimulator

lambda_true = 50.0

simulator = StochasticSimulator(
    "poisson", (lambda_true,),
    "lognormal", (10, 0.5),
    num_sim=100000,
    keep_all=True,
    seed=2025,
)

simulator.gen_agg_simulations()

events_per_year = simulator.all_simulations.groupby("year").size()
print("Theoretical mean:", lambda_true)
print("Simulated mean:",   events_per_year.mean())
\end{lstlisting}

\paragraph{Output interpretation.} The simulated mean claim count should
converge to $\lambda = 50$ as the number of simulations increases. Small
deviations are due to simulation noise and are bounded by the standard
error $\sqrt{\lambda / M}$, where $M$ is the number of simulations.

\paragraph{Actuarial notes.} This kind of internal consistency check
should be standard practice when validating any simulation-based model.
The validation chapter (Chapter~\ref{chap:validation}) discusses similar
checks for severity, aggregate losses, and claim development.
